%% cv.tex  –  Pingzhi Tang's CV
\documentclass[11pt,a4paper]{article}

% ── Packages ──────────────────────────────────────────────
\usepackage[margin=1in]{geometry}
\usepackage[T1]{fontenc}
\usepackage{lmodern}
\usepackage{fontspec}
\usepackage{enumitem}
\usepackage{xcolor}
\usepackage[colorlinks=true,urlcolor=linkblue,linkcolor=linkblue]{hyperref}
\usepackage{fancyhdr}
\usepackage{microtype}
\usepackage{textcomp}
\usepackage{bm}
\usepackage{fontawesome5}

% ── Colours ───────────────────────────────────────────────
\definecolor{headerblue}{RGB}{31,73,105}
\definecolor{linkblue}{RGB}{38,103,138}

% ── Page style / footer ──────────────────────────────────
\pagestyle{fancy}
\fancyhf{}
\renewcommand{\headrulewidth}{0pt}
\fancyfoot[C]{\footnotesize\itshape Pingzhi Tang~--~Page \thepage\ of \pageref*{LastPage}}

% ── Section font ────────────────────────────────────────
\newfontfamily\sectionfont{Gill Sans}
\renewcommand{\section}[1]{%
  \vspace{0.9em}%
  \noindent{\large\bfseries\sectionfont\color{headerblue}#1}%
  \vspace{-0.5em}%
  \par\noindent{\color{headerblue}\rule{\textwidth}{0.8pt}}%
  \vspace{0.3em}%
  \par
}

% ── List settings ────────────────────────────────────────
\setlist[itemize]{leftmargin=1.2em, itemsep=0pt, parsep=2pt, topsep=2pt}

% ── Helper: bold-italic name ─────────────────────────────
\newcommand{\myname}{\textbf{\textit{Pingzhi Tang}}}

% ══════════════════════════════════════════════════════════
\begin{document}

% ── Right-aligned "Last updated" ─────────────────────────
\begin{flushright}
  \vspace{-1em}
  \itshape Last updated: \today
\end{flushright}

\vspace{-0.5em}

% ── Name ─────────────────────────────────────────────────
\begin{center}
  {\Huge\bfseries\sectionfont\color{headerblue} Pingzhi Tang}\\[0.6em]
  \href{mailto:stanleytang@stu.pku.edu.cn}{\faEnvelope\ stanleytang@stu.pku.edu.cn}
  \quad \faPhone\ +86 182 5827 6580
  \quad \href{https://pingzhitang.cv}{\faGlobe\ pingzhitang.cv}\\[0.2em]
  \href{https://github.com/Stanleytowne}{\faGithub\ Stanleytowne}
  \quad \href{https://scholar.google.cz/citations?user=UYRV6y0AAAAJ}{\faGraduationCap\ Google Scholar}
\end{center}

% ── EDUCATION ────────────────────────────────────────────
\section{Education}

\noindent
\textbf{Peking University}, B.S.\ in Artificial Intelligence \hfill Sept 2023 -- present\\
General Artificial Intelligence Experimental Program (Tong Class)
\begin{itemize}
  \item \textit{GPA: 3.87/4.00; Major GPA: 3.92/4.00; \textbf{Rank 1/39}}
  \item Related courses: Linear Algebra, Probability and Statistics, Discrete Mathematics, Introduction to Computer Systems, Computer Networks, Machine Learning, Large Language Models and Alignment
\end{itemize}

% ── RESEARCH INTERESTS ───────────────────────────────────
\section{Research Interests}

\noindent
My research centers on optimizing the computational efficiency of Foundation Models with the ultimate goal of achieving accessible and scalable intelligence. I leverage system-level insights to innovate in model architecture and inference, viewing efficiency not merely as cost reduction, but as the key enabler for scaling reasoning and alignment.
\begin{itemize}
  \item \textbf{Inference Optimization:} KV Cache Compression, Speculative Decoding, Quantization, Model Pruning.
  \item \textbf{Next-Gen Architectures:} (Linear) Attention, Diffusion LLMs, SSMs.
  \item \textbf{Efficient and Fast Adaptation}: PEFT, efficient RL.
  \item \textbf{Reasoning \& RL:} Developing data-efficient RL algorithms and exploring scalable reasoning paradigms.
  \item \textbf{Continual Adaptation and Long-term Memory}: Enabling models to continuously learn and retain knowledge over extended contexts and tasks.
\end{itemize}

% ── PUBLICATIONS ─────────────────────────────────────────
\section{Publications}

\noindent
{[1]} \textbf{CLOVER: Cross-Layer Orthogonal Vectors Pruning and Fine-Tuning}\\
Fanxu Meng, \myname, Fan Jiang, Muhan Zhang\\
\href{https://arxiv.org/abs/2411.17426}{arxiv.org/abs/2411.17426}\quad (\textbf{ICML 2025})

\medskip\noindent
{[2]} \textbf{TransMLA: Multi-Head Latent Attention Is All You Need}\\
Fanxu Meng*, \myname*, Zengwei Yao, Xing Sun, Muhan Zhang (*equal contribution)\\
\href{https://arxiv.org/abs/2502.07864}{arxiv.org/abs/2502.07864}\quad (\textbf{NeurIPS 2025, spotlight})

\medskip\noindent
{[3]} \textbf{HD-PiSSA: High-Rank Distributed Orthogonal Adaptation}\\
Yiding Wang, Fanxu Meng, Xuefeng Zhang, Fan Jiang, \myname, Muhan Zhang\\
\href{https://arxiv.org/abs/2505.18777}{arxiv.org/abs/2505.18777}\quad (\textbf{EMNLP 2025, oral})

\medskip\noindent
{[4]} \textbf{TPLA: Tensor Parallel Latent Attention for Efficient Disaggregated Prefill Decode Inference}\\
Xiaojuan Tang*, Fanxu Meng*, \myname, Yuxuan Wang, Di Yin, Xing Sun, Muhan Zhang\\
\href{https://arxiv.org/abs/2508.15881}{arxiv.org/abs/2508.15881}\quad (\textbf{ASPLOS 2026})

\medskip\noindent
{[5]} \textbf{Orchestrating Dual-Boundaries: An Arithmetic Intensity Inspired Acceleration Framework for Diffusion Language Models}\\
Linye Wei, Wenjue Chen, \myname, Xiaotian Guo, Le Ye, Runsheng Wang, Meng Li\\
\href{https://arxiv.org/abs/2511.21759}{arxiv.org/abs/2511.21759}\quad (\textbf{DAC 2026})

\medskip\noindent
{[6]} \textbf{Knowledge is Not Enough: Injecting RL Skills for Continual Adaptation}\\
\myname*, Yiding Wang*, Muhan Zhang (*equal contribution)\\
\href{https://arxiv.org/abs/2601.11258}{arxiv.org/abs/2601.11258}\quad (\textbf{ACL 2026})

\medskip\noindent
{[7]} \textbf{Breaking the Blocks: Continuous Low-Rank Decomposed Scaling for Unified LLM Quantization and Adaptation}\\
\myname*, Ruijie Zhou*, Fanxu Meng, Wenjie Pei, Muhan Zhang (*equal contribution)\\
\href{https://arxiv.org/abs/2601.22716}{arxiv.org/abs/2601.22716}\quad (\textbf{Preprint})

\medskip\noindent
{[8]} \textbf{TEAM: Temporal-Spatial Consistency Guided Expert Activation for MoE Diffusion Language Model Acceleration}\\
Linye Wei, Zixiang Luo, \myname, Meng Li\\
\href{https://arxiv.org/abs/2602.08404}{arxiv.org/abs/2602.08404}\quad (\textbf{ICML 2026})

% ── RESEARCH EXPERIENCE ─────────────────────────────────
\section{Research Experience}

\noindent
\textbf{M$\bm{\mu}$ Lab}, Peking University \hfill July 2024 -- present\\
\textit{Supervisor: Prof.\ Muhan Zhang}
\begin{itemize}
  \item Developed a novel large model pruning and fine-tuning method using cross-layer singular value decomposition. Published at \textbf{ICML 2025} [\textbf{paper 1}].
  \item Created a practical method to transform standard multi-head attention into multi-head latent attention, achieving up to 11x inference acceleration with almost no loss. Published at \textbf{NeurIPS 2025 (Spotlight)} [\textbf{paper 2}].
  \item Invented High-rank Distributed PiSSA, a distributed PEFT method for LLMs that increases update flexibility by assigning different weight components to each GPU. Published at \textbf{EMNLP 2025} [\textbf{paper 3}].
  \item Proposed Parametric Skill Transfer (PaST), a modular framework that extracts domain-agnostic RL skill vectors and injects them into SFT-adapted models for efficient continual adaptation. Published at \textbf{ACL 2026} [\textbf{paper 6}].
  \item Proposed Low-Rank Decomposed Scaling (LoRDS), a unified framework for LLM quantization and adaptation that represents element-wise scaling as continuous low-rank matrices, enabling high-rank weight updates without additional inference cost [\textbf{paper 7}].
\end{itemize}

\medskip\noindent
\textbf{PKU-SEC-Lab}, Peking University \hfill Sept 2025 -- Feb 2026\\
\textit{Supervisor: Prof.\ Meng Li}
\begin{itemize}
  \item Developed an arithmetic intensity inspired acceleration framework for Diffusion Language Models, orchestrating compute- and memory-bound operations for efficient inference. Published at \textbf{DAC 2026} [\textbf{paper 5}].
  \item Proposed TEAM, a temporal-spatial consistency guided expert activation method for accelerating MoE Diffusion Language Models. Published at \textbf{ICML 2026} [\textbf{paper 8}].
\end{itemize}

% ── WORK EXPERIENCE ──────────────────────────────────────
\section{Work Experience}

\noindent
\textbf{Tencent}, Research Intern \hfill June 2025 -- Sept 2025
\begin{itemize}
  \item Co-invented a novel attention mechanism that addresses the KV cache overhead issue in MLA tensor parallelism. By utilizing tensor-parallel separation on top of MLA, this approach achieves a two-fold acceleration during inference. Published at \textbf{ASPLOS 2026} [\textbf{paper 4}].
  \item Attempted to integrate our TransMLA method [\textbf{paper 2}] into the training process of industry models.
\end{itemize}

% ── SKILLS ───────────────────────────────────────────────
\section{Skills}

\noindent
\textbf{Programming:} Proficient with Python, C++, Git and LaTeX

\medskip\noindent
\textbf{Mathematics:} Good understanding of calculus, linear algebra and probability

\medskip\noindent
\textbf{Languages:} English (fluent), Mandarin (native), French (beginner)

% ── ACADEMIC SERVICE ─────────────────────────────────────
\section{Academic Service}

\noindent
\textbf{Reviewer:} ICML 2026 (Gold Reviewer)

% ── ACHIEVEMENTS ─────────────────────────────────────────
\section{Achievements}

\noindent
\textbf{Taotian Scholarship, only 4 recipients annually} \hfill 2025\\
Institute for Artificial Intelligence, Peking University

\medskip\noindent
\textbf{Dean's Scholarship, top 2\%} \hfill 2025\\
Institute for Artificial Intelligence, Peking University

\medskip\noindent
\textbf{Soong Ching Ling ``Future Scholarship'', top 5\%} \hfill 2024\\
Institute for Artificial Intelligence, Peking University

\medskip\noindent
\textbf{Second-Class Scholarship, top 5\%} \hfill 2023 \& 2024\\
Peking University

\medskip\noindent
\textbf{``Outstanding Student'' Honor, top 10\%} \hfill 2023 \& 2024\\
Peking University

\label{LastPage}
\end{document}
